% Generated draft: requires mathematical and source review.
\documentclass[11pt]{article}
\usepackage[margin=25mm]{geometry}
\usepackage{fontspec}
\setmainfont{Latin Modern Roman}
\usepackage{amsmath,amssymb,mathtools,unicode-math}
\setmathfont{Latin Modern Math}
\usepackage{enumitem,needspace,xurl,hyperref}
\setlength{\parindent}{0pt}
\setlength{\parskip}{5pt}
\setlength{\emergencystretch}{3em}
\newcommand{\N}{\mathbb{N}}
\newcommand{\Z}{\mathbb{Z}}
\newcommand{\Q}{\mathbb{Q}}
\newcommand{\R}{\mathbb{R}}
\newcommand{\C}{\mathbb{C}}
\newcommand{\F}{\mathbb{F}}
\newcommand{\A}{\mathbb{A}}
\newcommand{\PP}{\mathbb{P}}
\newcommand{\E}{\mathbb{E}}
\newcommand{\eps}{\varepsilon}
\newcommand{\dd}{\,\mathrm{d}}
\begin{document}
\section*{UnsolvedMath ID: 30001908}
\textbf{Source problem code:} catalogue entry
\subsection{Definitions and mathematical statement}
The supplied question is
\[
\text{How can one compute }H_*(\Omega^\infty X)\text{ from knowledge of }H_*(X)\text{?}
\]
In the conventional spectrum interpretation of this notation, \(X\) is a spectrum and \(\Omega^\infty X\) is its derived zeroth infinite loop space. The record does not explicitly specify this domain, a class of admissible spectra, or homology coefficients.

For clarity, one possible interpretation, conditional on choosing a coefficient ring \(R\) and ordinary homology, uses
\[
H_n(X;R)=\pi_n(HR\wedge X),\qquad n\in\Z,
\]
where \(HR\) is the Eilenberg--Mac Lane spectrum of \(R\) and the smash product is derived, and
\[
H_n(\Omega^\infty X;R)=H_n^{\mathrm{sing}}(\Omega^\infty X;R),\qquad n\geq 0,
\]
where the right-hand side is singular homology. Under these explicitly conditional conventions, the question asks for a method to compute the latter graded homology from knowledge of the former.

The source does not specify whether the input consists only of graded homology groups or also includes operations or other structure. It also supplies no connectivity, component, localization, or finiteness hypotheses. Consequently, this is a computational research question with unresolved scope, rather than a fully specified assertion that such a method exists or that the output is determined by graded homology groups alone.
\subsection{Short English statement}
The problem asks how to recover the homology of the infinite loop space \(\Omega^\infty X\) from the homology of \(X\). A precise formulation requires specifying what kind of object \(X\) is, the homology theory and coefficients, the permitted hypotheses, and what structure is included in the input. The supplied record labels the problem open; that status is supplied dated curation.
\subsection{Sources}
\begin{itemize}
\item Local UnsolvedMath registry, permanent ID \(30001908\), title \emph{Homology of Infinite Loop Spaces}, source problem code \texttt{catalogue entry}, in the Oberwolfach Reports Open Problems collection. The record lists proposed year \(2011\) and supplies no named proposer.
\item Supplied report citation: \emph{Homotopy Theory (2012)}, Oberwolfach Reports, pp.\ 2623--2667, DOI \texttt{10.4171/owr/2011/46}. Supplied source URL and literature-source link, labelled ``original OWR report'': \url{https://doi.org/10.4171/owr/2011/46}. The background describes the report as documenting mathematical workshops and their open-problem discussions. This citation is reproduced from the record and has not been independently inspected.
\item Supplied literature assessment, checked on \texttt{2026-08-21}: the record labels the problem open and reports important special-case calculations and structural tools, while stating that no general procedure recovering \(H_*(\Omega^\infty X)\) from \(H_*(X)\) is known. This is supplied dated curation derived from the OWR material, not an independently verified literature assessment. The supplied background expressly calls for independent verification of time-sensitive status and cited claims.
\end{itemize}
\end{document}
